%-------------------------------------
% This template comes from Anish Athalye (Unofficial University of Cambridge Poster Template). 

% The Poster Template has been modified by Dr. Rahul Raoniar to fulfill B.Tech/Master/Ph.D./PostDoc student's poster presentation requirements.

% Description: I made this unofficial Poster Template for the Indian Institute of Technology Bombay (IITB). Feel free to use it, modify it, and share it. 

% Thank you note: A huge thanks goes to Anish Athalye for the original template.
%-------------------------------------

\documentclass[final]{beamer}

% ====================
% Packages
% ====================

\usepackage[T1]{fontenc}
\usepackage{lmodern}
\usepackage[size=custom,width=120,height=72,scale=1.0]{beamerposter}
\usetheme{gemini}
\usecolortheme{cam}
\usepackage{graphicx}
\usepackage{booktabs}
\usepackage{tikz}
\usepackage{pgfplots}
\pgfplotsset{compat=1.14}
\usepackage{anyfontsize}
\usepackage{xcolor}
\usepackage[skip=2pt,font=normalsize]{subcaption}
% ====================
% Lengths
% ====================

% If you have N columns, choose \sepwidth and \colwidth such that
% (N+1)*\sepwidth + N*\colwidth = \paperwidth
\newlength{\sepwidth}
\newlength{\colwidth}
\setlength{\sepwidth}{0.025\paperwidth}
\setlength{\colwidth}{0.3\paperwidth}

\newcommand{\separatorcolumn}{\begin{column}{\sepwidth}\end{column}}

% ====================
% Title
% ====================

\title{Poster Title, Poster Title, Poster Title, Poster Title,\\ Poster Title, Poster Title, Poster Title, Poster Title}

\author{Author1 \inst{1} \and Author2 \inst{2} \and Author3 \inst{2}}

\institute[shortinst]{\inst{1} Indian Institute of Technology Bombay \samelineand \inst{2} Indian Institute of Technology Guwahati}

% ====================
% Footer (optional)
% ====================

\footercontent{
  \href{https://www.lipsum.com}{\textbf{Your Portfolio Website}} \hfill
  \textbf{$14^{th}$ TPMDC, IIT Bombay, Mumbai, India} \hfill
  \href{mailto:your.email@provider-name}{\textbf{your.email@provider-name}}}
% (can be left out to remove footer)

% ====================
% Logo (optional)
% ====================

% use this to include logos on the left and/or right side of the header:
% \logoright{\includegraphics[height=7cm]{logo1.pdf}}
% \logoleft{\includegraphics[height=7cm]{logo2.pdf}}

% ====================
% Body
% ====================

\begin{document}

% Refer to https://github.com/k4rtik/uchicago-poster
% logo: https://www.cam.ac.uk/brand-resources/about-the-logo/logo-downloads
\addtobeamertemplate{headline}{}
{
    \begin{tikzpicture}[remember picture,overlay]
      \node [anchor=north west, inner sep=3cm] at ([xshift=5cm,yshift=1.5cm]current page.north west)
      {\includegraphics[height=9cm]{logos/iitb_logo.png}}; 
    \end{tikzpicture}
}

\addtobeamertemplate{headline}{}
{
    \begin{tikzpicture}[remember picture,overlay]
      \node [anchor=north west, inner sep=3cm] at ([xshift=100cm,yshift=1.5cm]current page.north west)
      {\includegraphics[height=8.5cm]{logos/tpmdc_logo.png}}; 
    \end{tikzpicture}
}

\begin{frame}[t]
\begin{columns}[t]
\separatorcolumn

\begin{column}{\colwidth}

% -------------------------------
% Section: Introduction
% -------------------------------

\begin{block}{Introduction}

    Lorem ipsum dolor sit amet, consectetur adipiscing elit. Ut blandit sed erat non vehicula. 

    Donec vel auctor ex. Donec vehicula a dui et imperdiet. Vivamus congue fringilla magna. Aenean iaculis metus ac pulvinar hendrerit.Donec vel auctor ex. Donec vehicula a dui et imperdiet. Vivamus congue fringilla magna. Aenean iaculis metus ac pulvinar hendrerit.Donec vel auctor ex. Donec vehicula a dui et imperdiet. Vivamus congue fringilla magna. Aenean iaculis metus ac pulvinar hendrerit.

    \heading{A heading inside a block}

    Integer elementum ante sed ex dapibus,  $u^4 + v^4$ sit amet laoreet ex iaculis. Duis accumsan commodo quam eu rutrum. Sed id risus eget justo luctus sollicitudin. Nulla faucibus lectus diam, vel pellentesque turpis fermentum ac.

\end{block}


% -------------------------------
% Section: Literature Review
% -------------------------------
\begin{block}{Literature Review}

    Duis sed lacinia augue. Sed suscipit nisl posuere enim eleifend, at cursus leo tincidunt. Sed blandit, metus a efficitur sagittis, sem quam dapibus elit, sed bibendum enim tellus et ligula \cite{li-1984}.

    A different kind of highlighted block.

    \begin{align*}
	x^2 - 3x -1 &= 0 \\
	x^3 + 4x + 3 &= 12 \\
	x - 3x &= 4 
    \end{align*}

    Duis sed lacinia augue. Sed suscipit nisl posuere enim eleifend, at cursus leo tincidunt.
    
    \begin{equation}
	\sqrt{x^2 + y^2} = z^2
    \end{equation}

    \heading{Research Gaps}

    Fusce viverra urna quis odio imperdiet, ut malesuada sem ultricies. Cras efficitur suscipit mauris eget viverra. Proin facilisis, nisl ac aliquam feugiat, velit erat volutpat lectus, at ultrices nisi mi ac erat. Vivamus vel molestie nibh, in bibendum nisl. Integer eu lacinia augue, nec blandit magna.

\end{block}

% -------------------------------
% Section: Research Objectives
% -------------------------------

\begin{alertblock}{Research Objectives}

   Enter your research objectives pointwise here.


    \begin{itemize}
      \item \textbf{Objective 1:} in sodales euismod, ligula elit feugiat velit, vel sodales sapien metus non massa.
      \item \textbf{Objective 2:} congue nunc vel commodo. Fusce efficitur egestas nibh sit amet porttitor. Mauris bibendum porttitor augue, at venenatis sapien fermentum in. 
      \item \textbf{Objective 3:} augue, nec blandit magna. Donec pharetra egestas urna at lobortis. 
    \end{itemize}

\end{alertblock}

\end{column}

% -------------------------------
% Section: Study Methodology
% -------------------------------
\separatorcolumn

\begin{column}{\colwidth}

  \begin{block}{Study Methodology}

    Mauris pretium vel erat quis accumsan. Ut pellentesque gravida tortor, vel vehicula ex.  

    \begin{enumerate}
      \item \textbf{Cras consequat congue} nunc vel commodo. Fusce efficitur egestas nibh sit amet porttitor.
      \item \textbf{Donec pharetra egestas} urna at lobortis. Phasellus sagittis, erat in elementum facilisis, urna augue ullamcorper orci, vel ultrices magna lorem sit amet arcu.
    \end{enumerate}


    Mauris eget commodo sapien, ut consequat mauris. Mauris in ligula luctus, tincidunt est eget, porttitor ante.

    % \begin{figure}
    %   \centering
    %   \begin{tikzpicture}
    %     \begin{axis}[
    %         scale only axis,
    %         no markers,
    %         domain=0:2*pi,
    %         samples=100,
    %         axis lines=center,
    %         axis line style={-},
    %         ticks=none]
    %       \addplot[red] {sin(deg(x))};
    %       \addplot[blue] {cos(deg(x))};
    %     \end{axis}
    %   \end{tikzpicture}
    %   \caption{Another figure caption.}
    % \end{figure}

\input{images/image1}

  \end{block}


% -------------------------------
% Section: Results and Discussion
% -------------------------------

\begin{block}{Results and Discussion}

    Cras consequat congue nunc vel commodo. Fusce efficitur egestas nibh sit amet porttitor. Mauris bibendum porttitor augue, at venenatis sapien fermentum in. Fusce viverra urna quis odio imperdiet, ut malesuada sem ultricies. Cras efficitur suscipit mauris eget viverra.

    \input{tables/table1}


\end{block}

% -------------------------------
% Section: Conclusions
% -------------------------------

\begin{block}{Conclusions}
    Write down the conclusions point-wise.
    \begin{itemize}
      \item \textbf{Sed nec} bibendum urna. Integer vulputate magna eget felis pellentesque, et porta sapien dictum.  
      \item \textbf{NAliquam dui leo}, tempus vel orci commodo, dignissim pulvinar velit. 
      \item \textbf{Nunc dapibus est} in nisi dictum, id tristique dolor lobortis.  
    \end{itemize}

  \end{block}

\end{column}

  % -------------------------------
  % Section: What is already known about this subject?
  % -------------------------------
\separatorcolumn

\begin{column}{\colwidth}

  \begin{exampleblock}{What is already known about this subject?}

    \begin{itemize}
      \item \textbf{Praesent non} cursus magna, ut imperdiet enim. 
      \item \textbf{Nunc dapibus} est in nisi dictum, id tristique dolor lobortis.
      \item \textbf{Aenean egestas} tincidunt est id interdum. 
    \end{itemize}

  \end{exampleblock}

  % -------------------------------
  % Section: What does this study add?
  % -------------------------------
  \begin{exampleblock}{What does this study add?}
    \begin{itemize}
      \item \textbf{Sed vel} luctus sapien, sit amet feugiat diam.
      \item \textbf{Phasellus maximus} purus vitae tempor bibendum. 
      \item \textbf{Nunc fermentum}, felis et mattis euismod, turpis dui tincidunt magna
    \end{itemize}

  \end{exampleblock}

  \begin{block}{References}

    % \nocite{*}
    \footnotesize{\bibliographystyle{plain}\bibliography{poster}}

  \end{block}

\end{column}

\separatorcolumn
\end{columns}
\end{frame}

\end{document}
